\documentclass[11pt,a4paper]{article}
% Společný preambule pro cvičnou písemku a klíč řešení (XeLaTeX)
\usepackage[a4paper,left=2.1cm,right=2.1cm,top=2.4cm,bottom=2.2cm,headheight=14pt,footskip=1.1cm]{geometry}
\usepackage{amsmath}
\usepackage{fontspec}
\usepackage{unicode-math}
\usepackage{polyglossia}
\setdefaultlanguage{czech}
\setmainfont{Latin Modern Roman}
\setmathfont{Latin Modern Math}
\setsansfont{Inter}[Scale=0.92]
\usepackage{microtype}
\usepackage{xcolor}
\usepackage{enumitem}
\usepackage{array}
\usepackage{multicol}
\usepackage{needspace}
\usepackage{xspace}
\usepackage{fancyhdr}
\usepackage{lastpage}
\usepackage{tikz}
\usetikzlibrary{calc}
\usepackage[skins,breakable]{tcolorbox}

\definecolor{ink}{HTML}{1F2937}
\definecolor{accent}{HTML}{2B4C7E}
\definecolor{accentlight}{HTML}{EAF0FA}
\definecolor{muted}{HTML}{6B7280}
\definecolor{rulegray}{HTML}{B8C2D0}
\definecolor{rulelight}{HTML}{E1E6EE}
\definecolor{good}{HTML}{1F6F4A}
\definecolor{goodlight}{HTML}{EAF6EF}
\definecolor{warnlight}{HTML}{FFF6E5}
\definecolor{warn}{HTML}{9A5B00}

\color{ink}
\linespread{1.12}
\setlength{\parindent}{0pt}
\setlength{\parskip}{0pt}
\raggedright
\setlist{nosep}

% Záhlaví a zápatí
\pagestyle{fancy}
\fancyhf{}
\renewcommand{\headrulewidth}{0pt}
\renewcommand{\footrulewidth}{0pt}
\fancyhead[L]{\sffamily\footnotesize\color{muted}Matematické dovednosti · Množiny}
\fancyfoot[C]{\sffamily\footnotesize\color{muted}\thepage\,/\,\pageref{LastPage}}

% Pruh s názvem části
\newcommand{\cast}[2]{%
  \par\Needspace*{11cm}\vspace{1.6em}\noindent
  \begin{tcolorbox}[enhanced,colback=accent,colframe=accent,boxrule=0pt,arc=3pt,left=10pt,right=10pt,top=5pt,bottom=5pt,before skip=0pt,after skip=6pt]
    {\sffamily\bfseries\color{white}\small Část #1\enspace\textcolor{white!70!accent}{|}\enspace #2}
  \end{tcolorbox}\par\nopagebreak}

% Úloha: číslo, text
\newcommand{\uloha}[1]{\par\vspace{1.15em}\noindent\textcolor{accent}{\sffamily\bfseries #1}\enspace}
% Body
\newcommand{\bodu}[1]{\hfill\mbox{\sffamily\footnotesize\color{muted}#1\,b}}

% Místo na řešení: linkovaný box dané výšky
\newcommand{\misto}[1]{%
  \par\vspace{0.35em}
  \begin{tcolorbox}[enhanced,height=#1,colback=white,colframe=rulegray,boxrule=0.5pt,arc=3pt,left=0pt,right=0pt,top=0pt,bottom=0pt,before skip=0pt,after skip=0pt,
    underlay={\begin{scope}\clip (frame.south west) rectangle (frame.north east);
      \foreach \i in {1,...,40}{\draw[rulelight,line width=0.4pt] ([yshift=-\i*0.9cm]frame.north west) -- ([yshift=-\i*0.9cm]frame.north east);}
    \end{scope}}]
  \end{tcolorbox}\par}

% Rámeček s poznámkou
\newtcolorbox{tip}[1][]{enhanced,breakable,colback=accentlight,colframe=accentlight,boxrule=0pt,arc=4pt,left=10pt,right=10pt,top=7pt,bottom=7pt,before skip=8pt,after skip=8pt,#1}
\newtcolorbox{varovani}[1][]{enhanced,breakable,colback=warnlight,colframe=warnlight,boxrule=0pt,arc=4pt,left=10pt,right=10pt,top=7pt,bottom=7pt,before skip=8pt,after skip=8pt,#1}

% Řešení (pro klíč)
\newtcolorbox{reseni}{enhanced,breakable,colback=goodlight,colframe=goodlight,boxrule=0pt,arc=3pt,left=10pt,right=8pt,top=6pt,bottom=6pt,before skip=5pt,after skip=2pt,
  borderline west={2.2pt}{0pt}{good}}
\newcommand{\pravda}{\textcolor{good}{\sffamily\bfseries pravda}\xspace}
\newcommand{\nepravda}{\textcolor{warn}{\sffamily\bfseries nepravda}\xspace}

% Zkratky
\newcommand{\N}{\mathbb N}
\newcommand{\Z}{\mathbb Z}
\newcommand{\Pot}{\mathcal P}
\newcommand{\pr}{\mathrel{\mid}}
\newcommand{\nepr}{\mathrel{\nmid}}
\newcommand{\sm}{\setminus}

% Krátké položky ve dvou sloupcích / jeden sloupec
\newenvironment{dvasl}{\begin{multicols}{2}\begin{enumerate}[label=(\alph*),leftmargin=2.2em,itemsep=3pt,topsep=3pt]}{\end{enumerate}\end{multicols}}
\newenvironment{pol}{\begin{enumerate}[label=(\alph*),leftmargin=2.2em,itemsep=3pt,topsep=4pt]}{\end{enumerate}}
\setlength{\columnsep}{1.2cm}

\newcommand{\blacksquare}{\mathord{\rule{0.85ex}{0.85ex}}}

\fancyhead[R]{\sffamily\footnotesize\color{muted}Řešení cvičné písemky}

\newcommand{\zad}[2]{\par\Needspace*{4cm}\vspace{1.1em}\noindent\textcolor{accent}{\sffamily\bfseries #1}\enspace #2\par}

\begin{document}

{\sffamily\bfseries\fontsize{25}{28}\selectfont\color{accent} Množiny}\par
\vspace{2pt}
{\sffamily\large\color{muted} Řešení cvičné písemky}\par
\vspace{10pt}
{\color{rulegray}\rule{\linewidth}{0.6pt}}\par
\vspace{8pt}

\begin{varovani}
{\sffamily\bfseries\small Nejdřív řeš, potom čti.}\par\vspace{2pt}
Řešení má smysl až po vlastním pokusu. U důkazů je podstatná \emph{stavba} argumentu (co předpokládám, co dokazuji, který krok je ekvivalence). Stejné závěry z jiného postupu jsou také správně.
\end{varovani}

% ==================== ČÁST 1 ====================
\cast{1}{Zápisy množin}

\zad{1.1}{Výčtem prvků.}
\begin{reseni}
(a) $A=\{-4,-3,-2,-1,0,1,2,3,4\}$ (pro $|x|=4$ je $16<20$, pro $|x|=5$ je $25\not<20$; pozor na záporná čísla a nulu).\\
(b) $B=\{4,6,8,12,24\}$ (dělitelé 24 jsou $1,2,3,4,6,8,12,24$, ponecháme ty větší než 3).\\
(c) $C=\{2,5,8,11\}$.\\
(d) $D=\{3,6,9,12,15,18\}$.
\end{reseni}

\zad{1.2}{Množinovým zápisem.}
\begin{reseni}
(a) $\{5k\mid k\in\N\}=\{n\in\N\mid 5\pr n\}=\{n\in\N\mid\exists k\in\N:\ n=5k\}$.\\
(b) $\{k^2\mid k\in\N\}=\{n\in\N\mid\exists k\in\N:\ n=k^2\}$.\\
(c) $\{3k+2\mid k\in\N_0\}=\{n\in\N\mid\exists k\in\N_0:\ n=3k+2\}=\{n\in\N\mid 3\pr(n-2)\}$.\\
(d) $\{d\in\N\mid d\pr 30\}=\{1,2,3,5,6,10,15,30\}$.

\smallskip
{\sffamily\footnotesize\color{muted}Převod mezi zápisy: $\{f(k)\mid k\in K\}=\{y\mid\exists k\in K:\ y=f(k)\}$. Pozor, odkud $k$ běží ($\N$ vs.\ $\N_0$).}
\end{reseni}

\zad{1.3}{Zápis pomocí $X,Y\subseteq\N$.}
\begin{reseni}
(a) $\{x\in X\mid\forall y\in Y:\ y\pr x\}$\\
(b) $\{y\in Y\mid\exists x\in X:\ y<x\}$\\
(c) $\{x\in X\mid\forall y\in Y:\ x>y\}$\\
(d) $\{x_1+x_2\mid x_1,x_2\in X\}=\{n\in\N\mid\exists x_1,x_2\in X:\ n=x_1+x_2\}$\\
(e) $\{y\in Y\mid\exists x\in X:\ y=x^2\}$

\smallskip
{\sffamily\footnotesize\color{muted}„Každý/všechny“ $\to\forall$, „některý/nějaký“ $\to\exists$. Prvky výsledku jsou z $X$ (resp.\ $Y$), proto se vybírá z té množiny.}
\end{reseni}

\zad{1.4}{Právě jedna z množin $A,B$ (symetrický rozdíl).}
\begin{reseni}
(i) $(A\cup B)\setminus(A\cap B)$, nebo $(A\setminus B)\cup(B\setminus A)$.\\
(ii) $\{x\mid(x\in A\land x\notin B)\lor(x\in B\land x\notin A)\}$.
\end{reseni}

\zad{1.5}{$M=\{n\in\N\mid\forall k\in\N:\ (k\pr n\Rightarrow k=1\lor k=n)\}$.}
\begin{reseni}
Podmínka říká: jediní možní dělitelé $n$ jsou $1$ a $n$. To splňují \textbf{prvočísla} \emph{a také číslo $1$} (jeho jediný dělitel je $1$, tedy $k=1$ vždy platí). $M=\{1\}\cup\{\text{prvočísla}\}=\{1,2,3,5,7,11,\dots\}$.
\end{reseni}

% ==================== ČÁST 2 ====================
\cast{2}{Prvky, prázdná množina, počet prvků}

\zad{2.1}{Pravda / nepravda.}
\begin{reseni}
(a) \pravda: $\emptyset\subseteq A$ platí pro každou množinu $A$.\\
(b) \nepravda: $\{\emptyset\}$ má prvek $\emptyset$, který neleží v $\emptyset$.\\
(c) \nepravda: jediný prvek množiny $\{\{\emptyset\}\}$ je $\{\emptyset\}$, a to není $\emptyset$.\\
(d) \pravda: $\{\emptyset\}$ je přesně ten jediný prvek.\\
(e) \pravda: $\{1,2\}$ je třetím prvkem.\\
(f) \pravda: $1\in S$ a $2\in S$, tedy každý prvek $\{1,2\}$ leží v $S$.\\
(g) \pravda: $\{1\}$ je prvkem $\{1,\{1\}\}$, a $\{\{1\}\}$ obsahuje jen $\{1\}$.\\
(h) \pravda: $\Pot(\emptyset)=\{\emptyset\}$ (jedině $\emptyset$ je podmnožinou $\emptyset$).
\end{reseni}

\zad{2.2}{Počet prvků.}
\begin{reseni}
(a) $\{1,\{1\},\{1,1\},\{\{1\}\},\emptyset,\{\emptyset\},\{\}\}$: platí $\{1,1\}=\{1\}$ a $\{\}=\emptyset$. Různé prvky: $1$, $\{1\}$, $\{\{1\}\}$, $\emptyset$, $\{\emptyset\}$. \textbf{5 prvků.}\\
(b) $\{a,b,c,\{a,b\}\}$: \textbf{4 prvky.}\\
(c) $\{\emptyset,\{\emptyset,\emptyset\},\{\{\emptyset\}\}\}$, kde $\{\emptyset,\emptyset\}=\{\emptyset\}$. Prvky: $\emptyset$, $\{\emptyset\}$, $\{\{\emptyset\}\}$. \textbf{3 prvky.}\\
(d) $\Pot(\{a,b\})=\{\emptyset,\{a\},\{b\},\{a,b\}\}$, sjednocením s $\{a,b\}$ přibudou prvky $a,b$ (to jsou jiné objekty než podmnožiny). Celkem $\emptyset,\{a\},\{b\},\{a,b\},a,b$: \textbf{6 prvků.}
\end{reseni}

\zad{2.3}{Tabulka pro $S=\{1,2,\{3\},\{1,2\}\}$.}
\begin{reseni}
\renewcommand{\arraystretch}{1.25}
\begin{tabular}{@{}lcc@{}}
$x$ & $x\in S$ & $x\subseteq S$\\
$\{1\}$ & ne & \textbf{ano} ($1\in S$)\\
$\{3\}$ & \textbf{ano} & ne ($3\notin S$)\\
$\{1,2\}$ & \textbf{ano} & \textbf{ano} ($1,2\in S$)\\
$\{\{3\}\}$ & ne & \textbf{ano} ($\{3\}\in S$)\\
$\emptyset$ & ne & \textbf{ano}
\end{tabular}

\smallskip
{\sffamily\footnotesize\color{muted}$x\in S$: je $x$ \emph{přímo} prvkem seznamu? $x\subseteq S$: leží \emph{každý prvek} množiny $x$ v $S$?}
\end{reseni}

\zad{2.4}{Potenční množiny.}
\begin{reseni}
(a) $\Pot(\{1,\{2\}\})=\{\emptyset,\ \{1\},\ \{\{2\}\},\ \{1,\{2\}\}\}$ (4 prvky; množina má dva prvky: $1$ a $\{2\}$).\\
(b) $\Pot(\emptyset)=\{\emptyset\}$, tedy $\Pot(\Pot(\emptyset))=\Pot(\{\emptyset\})=\{\emptyset,\{\emptyset\}\}$.\\
(c) $\Pot(\{a\})=\{\emptyset,\{a\}\}$ má 2 prvky, proto $|\Pot(\Pot(\{a\}))|=2^2=4$.
\end{reseni}

% ==================== ČÁST 3 ====================
\cast{3}{Podmnožiny}

\zad{3.1}{$A=\{1,..,5\}$, $B=\{2,4\}$, $C=\{1,2,3\}$, $D=\{2,3,6\}$.}
\begin{reseni}
(a) \pravda\\
(b) \nepravda ($6\notin A$)\\
(c) \nepravda ($4\notin C$)\\
(d) \pravda (průnik je vždy podmnožinou obou)\\
(e) \nepravda ($C\cap D=\{2,3\}$ a $3\notin B$)\\
(f) \pravda (rozdíl je vždy podmnožinou prvního).
\end{reseni}

\zad{3.2}{Kvantifikátory.}
\begin{reseni}
(a) $\forall x\in A:\ x\in B$\\
(b) $\exists x\in A:\ x\notin B$\\
(c) $\forall x\in A:\ x\notin B$ (ekvivalentně $\neg\exists x\,(x\in A\land x\in B)$)\\
(d) $(\forall x\in A:\ x\in B)\land(\exists y\in B:\ y\notin A)$\\
(e) $(\exists x\in A:\ x\notin B)\lor(\exists y\in B:\ y\notin A)$, tj.\ $A\nsubseteq B$ nebo $B\nsubseteq A$.
\end{reseni}

\zad{3.3}{Dělitelnost a množiny.}
\begin{reseni}
(a) $A\subseteq B$ (každý násobek 6 je sudý), navíc $A\neq B$ (např.\ $2\in B\setminus A$).\\
(b) $A\subseteq C$ (každý násobek 6 je násobkem 3).\\
(c) \textbf{neporovnatelné}: $2\in B\setminus C$ a $3\in C\setminus B$.\\
(d) $A=B\cap C$: $n\in B\cap C\Leftrightarrow 2\pr n\land 3\pr n\Leftrightarrow 6\pr n$, protože 2 a 3 jsou nesoudělná.
\end{reseni}

\zad{3.4}{Důkazy.}
\begin{reseni}
\textbf{(a)} Nechť $A\subseteq B$. Vezměme libovolné $x\in A\cap C$. Pak $x\in A$ a $x\in C$. Z $A\subseteq B$ plyne $x\in B$. Tedy $x\in B$ a $x\in C$, takže $x\in B\cap C$. Protože $x$ bylo libovolné, $A\cap C\subseteq B\cap C$. \hfill$\blacksquare$

\smallskip
\textbf{(b)} Nechť $A\cup B=A\cap B$. \emph{$A\subseteq B$:} je-li $x\in A$, pak $x\in A\cup B=A\cap B$, tedy $x\in B$. \emph{$B\subseteq A$:} je-li $x\in B$, pak $x\in A\cup B=A\cap B$, tedy $x\in A$. Z obou inkluzí plyne $A=B$. \hfill$\blacksquare$
\end{reseni}

% ==================== ČÁST 4 ====================
\cast{4}{Sjednocení, průnik, rozdíl}

\zad{4.1}{$A=\{1,..,6\}$, $B=\{2,4,6,8\}$, $C=\{3,6,9\}$.}
\begin{reseni}
(a) $A\cup B=\{1,2,3,4,5,6,8\}$\\
(b) $A\cap B=\{2,4,6\}$\\
(c) $A\setminus B=\{1,3,5\}$\\
(d) $B\setminus A=\{8\}$\\
(e) $(A\cap B)\cup C=\{2,3,4,6,9\}$\\
(f) $A\setminus(B\cup C)$: $B\cup C=\{2,3,4,6,8,9\}$, tedy $\{1,5\}$\\
(g) $(A\cup C)\setminus B$: $A\cup C=\{1,\dots,6,9\}$, po odebrání $B$ zůstane $\{1,3,5,9\}$\\
(h) $A\cup B\cup C=\{1,2,3,4,5,6,8,9\}$, tedy $|A\cup B\cup C|=8$.
\end{reseni}

\zad{4.2}{$A$ sudá, $B$ násobky 3, $C=\{1,\dots,12\}$.}
\begin{reseni}
(a) $A\cap B\cap C=\{6,12\}$\\
(b) $(A\cup B)\cap C=\{2,3,4,6,8,9,10,12\}$\\
(c) $C\setminus(A\cup B)=\{1,5,7,11\}$\\
(d) $(A\setminus B)\cap C=\{2,4,8,10\}$
\end{reseni}

\zad{4.3}{Skupiny rovných množin.}
\begin{reseni}
Skupiny: \textbf{(a)=(b)}, \textbf{(c)=(d)}, \textbf{(e)=(f)}.

\smallskip
\emph{Důkaz (e)=(f):} pro libovolné $x$ platí
\[
\begin{aligned}
x\in(A\cup B)\setminus C&\Leftrightarrow(x\in A\lor x\in B)\land x\notin C\\
&\Leftrightarrow(x\in A\land x\notin C)\lor(x\in B\land x\notin C)\\
&\Leftrightarrow x\in(A\setminus C)\cup(B\setminus C).
\end{aligned}
\]
(Druhá ekvivalence je distributivita $\land$ přes $\lor$.)

\smallskip
\emph{Proč (a)=(b):} $x\in A\setminus(B\cap C)\Leftrightarrow x\in A\land\neg(x\in B\land x\in C)\Leftrightarrow x\in A\land(x\notin B\lor x\notin C)\Leftrightarrow(A\setminus B)\cup(A\setminus C)$. \emph{(c)=(d):} $x\in A\land x\notin B\land x\notin C$ vyjde v obou případech.

\smallskip
\emph{Protipříklad, že (a)$\ne$(c):} $A=B=\{1\}$, $C=\emptyset$. Pak $A\setminus(B\cap C)=\{1\}\setminus\emptyset=\{1\}$, ale $(A\setminus B)\cap(A\setminus C)=\emptyset\cap\{1\}=\emptyset$.
\end{reseni}

\zad{4.4}{Dokázat, nebo vyvrátit.}
\begin{reseni}
(a) \nepravda: $A=\{1\}$, $B=\{2\}$. $(A\setminus B)\cup B=\{1,2\}\neq\{1\}$. (Obecně vychází $A\cup B$.)\\[3pt]
(b) \pravda: $x\in A\cap(B\setminus C)\Leftrightarrow x\in A\land x\in B\land x\notin C\Leftrightarrow x\in(A\cap B)\setminus C$.\\[3pt]
(c) \pravda: $x\in A\setminus(A\setminus B)\Leftrightarrow x\in A\land\neg(x\in A\land x\notin B)\Leftrightarrow x\in A\land(x\notin A\lor x\in B)\Leftrightarrow x\in A\land x\in B$ (první část $x\in A\land x\notin A$ je spor).\\[3pt]
(d) \nepravda: $A=B=C=\{1\}$. $A\setminus(B\setminus C)=\{1\}\setminus\emptyset=\{1\}$, ale $(A\setminus B)\setminus C=\emptyset$. Rozdíl množin není asociativní.
\end{reseni}

\zad{4.5}{Velikosti.}
\begin{reseni}
(a) $|A\cup B|=12+9-4=17$, $|A\setminus B|=12-4=8$, $|A\times B|=12\cdot9=108$, $|\Pot(A\cap B)|=2^4=16$.\\
(b) Aspoň jeden sport hraje $18+14-7=25$ studentů, žádný tedy $30-25=5$.
\end{reseni}

% ==================== ČÁST 5 ====================
\cast{5}{Potenční množina, indexovaná sjednocení a průniky}

\zad{5.1}{$A=\{1,2\}$, tedy $\Pot(A)=\{\emptyset,\{1\},\{2\},\{1,2\}\}$.}
\begin{reseni}
(a) \pravda\\
(b) \nepravda ($1$ není podmnožina, ale prvek $A$)\\
(c) \pravda\\
(d) \pravda ($\{1\}\in\Pot(A)$, a $\{\{1\}\}$ obsahuje jen $\{1\}$)\\
(e) \nepravda ($1\in A$, ale $1\notin\Pot(A)$)\\
(f) \nepravda ($3\notin A$).
\end{reseni}

\zad{5.2}{$M_i=\{j\in\N\mid 2j\le i\}$, $i=1,\dots,5$.}
\begin{reseni}
$M_1=\emptyset$, $M_2=M_3=\{1\}$, $M_4=M_5=\{1,2\}$.

(a) $\bigcup M_i=\{1,2\}$\\
(b) $\bigcap M_i=\emptyset$ (kvůli $M_1=\emptyset$)\\
(c) $\bigcup\{M_i\}=\{\emptyset,\{1\},\{1,2\}\}$ (sbírá se samotné množiny $M_i$ jako prvky)\\
(d) $\bigcap\{M_i\}=\emptyset$ (jednoprvkové množiny $\{\emptyset\}$, $\{\{1\}\}$, $\{\{1,2\}\}$ nemají společný prvek)\\
(e) $\{\bigcup M_i\}=\{\{1,2\}\}$ (výsledek z (a) obalený závorkami)
\end{reseni}

\zad{5.3}{Potence a operace.}
\begin{reseni}
\textbf{(a)} Pro libovolnou množinu $X$:
\[
\begin{aligned}
X\in\Pot(A\cap B)&\Leftrightarrow X\subseteq A\cap B\Leftrightarrow(X\subseteq A\land X\subseteq B)\\
&\Leftrightarrow(X\in\Pot(A)\land X\in\Pot(B))\Leftrightarrow X\in\Pot(A)\cap\Pot(B).\ \blacksquare
\end{aligned}
\]
\textbf{(b)} $A=\{1\}$, $B=\{2\}$: $\{1,2\}\in\Pot(A\cup B)$, ale $\{1,2\}\notin\Pot(A)\cup\Pot(B)$, protože není podmnožinou $A$ ani $B$.

\textbf{(c)} Ano. Je-li $X\in\Pot(A)\cup\Pot(B)$, pak $X\subseteq A$ nebo $X\subseteq B$. Obojí je podmnožina $A\cup B$, tedy $X\in\Pot(A\cup B)$. Opačná inkluze obecně neplatí podle (b). \hfill$\blacksquare$
\end{reseni}

\zad{5.4}{Počty.}
\begin{reseni}
(a) $2^5=32$.\\
(b) $16$. Podmnožin se sudým počtem prvků je $\binom50+\binom52+\binom54=1+10+5=16$ (i $\emptyset$ má sudý počet 0).\\
(c) $\binom52=10$.\\
(d) $|A\times B|=6$, tedy $|\Pot(A\times B)|=2^6=64$.
\end{reseni}

% ==================== ČÁST 6 ====================
\newpage
\cast{6}{Cvičná písemka na čas: řešení s bodováním}

\zad{Úloha 1\ \ (4 b)}{Zápisy, 1~bod za každou správnou položku.}
\begin{reseni}
(a) $\{n\in\N\mid 4\pr n\lor 6\pr n\}$\\
(b) $\{x\in X\mid\forall y\in Y:\ x<y\}$\\
(c) $\{3x\mid x\in X\}=\{n\in\N\mid\exists x\in X:\ n=3x\}$\\
(d) $\{S\in\Pot(M)\mid |S|\ge2\}$ (nebo $\{S\mid S\subseteq M\land|S|\ge2\}$)
\end{reseni}

\zad{Úloha 2\ \ (4 b)}{Pravda / nepravda, 1~bod za položku se zdůvodněním.}
\begin{reseni}
(a) \pravda: $\emptyset\in\Pot(\{1\})$, takže jediný prvek $\{\emptyset\}$ leží v $\Pot(\{1\})$.\\
(b) \nepravda: jediný prvek $\{\{\emptyset\}\}$ je $\{\emptyset\}$, nikoli $\emptyset$.\\
(c) \pravda: $\{1\}\in\{1,\{1\}\}$.\\
(d) \nepravda: prvky $\Pot(\{1,2\})$ jsou \emph{podmnožiny} $\{1,2\}$, tedy $\emptyset,\{1\},\{2\},\{1,2\}$. Číslo $2$ je prvek $\{1,2\}$, ne jeho podmnožina; správně je $\{2\}\in\Pot(\{1,2\})$.
\end{reseni}

\zad{Úloha 3\ \ (4 b)}{Počty prvků.}
\begin{reseni}
(a) \textbf{2 b.} $\{a,\{a\},\{a,a\},\emptyset,\{\emptyset\},\{\{a\}\}\}$: $\{a,a\}=\{a\}$. Prvky: $a$, $\{a\}$, $\emptyset$, $\{\emptyset\}$, $\{\{a\}\}$ $\to$ \textbf{5 prvků} (1~b za počet, 1~b za výčet).\\
(b) \textbf{1 b.} $|\Pot(\{1\})|=2$, $|\Pot(\Pot(\{1\}))|=2^2=4$.\\
(c) \textbf{1 b.} $|A\times B|=3\cdot2=6$, $|A\cup B|=3+2-1=4$.
\end{reseni}

\zad{Úloha 4\ \ (4 b)}{Kvantifikátory a negace.}
\begin{reseni}
\textbf{2 b:} $\forall a\in A\ \exists b\in B:\ a<b$.\\
\textbf{2 b:} negace $\exists a\in A\ \forall b\in B:\ a\ge b$, slovně: „v $A$ existuje číslo, které je větší nebo rovno všem číslům z $B$“.

\smallskip
{\sffamily\footnotesize\color{muted}Postup: negace převrátí \emph{oba} kvantifikátory (zachová pořadí) a znegovanou vlastnost $a<b$ změní na $a\ge b$. Za neúplnou negaci (např.\ jen první kvantifikátor, nebo $a>b$) se strhává.}
\end{reseni}

\zad{Úloha 5\ \ (4 b)}{Rovnosti množin.}
\begin{reseni}
\textbf{(a) (2 b)} Pro libovolné $x$:
$x\in(A\cup B)\setminus C\Leftrightarrow(x\in A\lor x\in B)\land x\notin C\Leftrightarrow(x\in A\land x\notin C)\lor(x\in B\land x\notin C)\Leftrightarrow x\in(A\setminus C)\cup(B\setminus C)$. \hfill$\blacksquare$

\textbf{(b) (2 b)} Protipříklad: $A=B=C=\{1\}$. Pak $A\setminus(B\setminus C)=\{1\}\setminus\emptyset=\{1\}$, ale $(A\setminus B)\setminus C=\emptyset\setminus\{1\}=\emptyset$. Rovnost tedy neplatí. (Stačí \emph{jedna} konkrétní trojice množin, ale musí ověřit obě strany.)
\end{reseni}

\vspace{1.2em}
\begin{tip}
{\sffamily\bfseries\small Orientační hodnocení}\par\vspace{2pt}
\textbf{18–20 b} výborně, \textbf{14–17 b} dobře (zopakuj, co ti nešlo), \textbf{pod 14 b} projdi poznámky \textit{03–05} a zkus část 6 znovu za pár dní.
\end{tip}

\end{document}
